\chapter*{Annexes}
\addcontentsline{toc}{chapter}{Annexes}
\markboth{Annexes}{}
\stepcounter{chapter}
\addtocontents{lot}{\vspace{3.8mm}}
\addtocontents{lof}{\vspace{3.8mm}}

%===================== ANNEXE 1 =====================%
\section*{Annexe 1.~Liste complète des points d'entrée de l'API}
\addcontentsline{toc}{section}{Annexe 1.~Liste complète des points d'entrée de l'API}

\addcontentsline{lot}{table}{Annexe 1.1~~~Points d'entrée exposés par le micro-service ML}

Le tableau annexe~1.1 recense l'ensemble des points d'entrée REST exposés par le micro-service de Machine Learning (FastAPI).

{\raggedright \textbf{Tableau annexe 1.1~:}~Points d'entrée exposés par le micro-service ML}
\begin{longtable}[c]{
    | p{.30\textwidth}
    | p{.12\textwidth}
    | p{.50\textwidth} |
}
    \hline
    \textbf{Route} & \textbf{Méthode} & \textbf{Description} \\ \hline
    \endhead
    \texttt{/api/predict/carbon} & POST & Prédiction d'empreinte carbone -- transport \\ \hline
    \texttt{/api/predict/energy} & POST & Prédiction d'empreinte carbone -- énergie domestique \\ \hline
    \texttt{/api/predict/electronics} & POST & Prédiction d'empreinte carbone -- électronique \\ \hline
    \texttt{/api/predict/water} & POST & Prédiction de consommation d'eau \\ \hline
    \texttt{/api/predict/food} & POST & Prédiction d'empreinte carbone -- alimentation \\ \hline
    \texttt{/api/predict/waste} & POST & Prédiction d'empreinte carbone -- déchets \\ \hline
    \texttt{/api/forecast} & POST & Prévision de tendance à partir d'un historique \\ \hline
    \texttt{/api/models/metrics} & GET & Métriques complètes (MAE/RMSE/R\textsuperscript{2}) des 5 algorithmes, 6 catégories \\ \hline
    \texttt{/api/models/\{category\}/distribution} & GET & Percentiles de la distribution cible d'une catégorie \\ \hline
    \texttt{/health} & GET & Point de contrôle de santé du service (Docker healthcheck) \\ \hline
\end{longtable}

\addcontentsline{lot}{table}{Annexe 1.2~~~Exemple de requête et de réponse}

Le tableau annexe~1.2 illustre un exemple concret d'appel à l'API de prédiction pour la catégorie Transport.

{\raggedright \textbf{Tableau annexe 1.2~:}~Exemple de requête / réponse -- \texttt{POST /api/predict/carbon}}

\begin{lstlisting}[caption={Corps de la requête (JSON)}]
{
  "vehicleType": "car",
  "fuelType": "petrol",
  "distance": 20,
  "fuelConsumption": 6,
  "passengers": 1
}
\end{lstlisting}

\begin{lstlisting}[caption={Réponse renvoyée (JSON)}]
{
  "category": "carbon",
  "prediction": 2.77,
  "target": "co2e_kg",
  "unit": "kg CO2e",
  "greenScore": 95,
  "confidence": 0.982,
  "modelUsed": "gradient_boosting",
  "percentileRank": 5.0,
  "featureContributions": {
    "distance": 0.43,
    "fuelType": 0.27,
    "fuelConsumption": 0.19,
    "vehicleType": 0.08,
    "passengers": 0.03
  }
}
\end{lstlisting}

\newpage
%===================== ANNEXE 2 =====================%
\section*{Annexe 2.~Extrait de la configuration de déploiement Docker}
\addcontentsline{toc}{section}{Annexe 2.~Extrait de la configuration de déploiement Docker}

L'extrait ci-dessous, tiré du fichier \texttt{docker-compose.yml} du projet, illustre l'orchestration des deux conteneurs applicatifs et leur configuration de supervision (\textit{healthcheck}).

\begin{lstlisting}[caption={Extrait de docker-compose.yml}]
services:
  ml-service:
    build: ./ml-service
    ports:
      - "8000:8000"
    healthcheck:
      test: ["CMD", "python", "-c",
             "import urllib.request; urllib.request.urlopen('http://localhost:8000/health')"]
      interval: 30s
      timeout: 5s
      retries: 3
    restart: unless-stopped

  frontend:
    build:
      context: .
      args:
        VITE_SUPABASE_URL: "${VITE_SUPABASE_URL}"
        VITE_ML_API_URL: "${VITE_ML_API_URL:-http://localhost:8000}"
    ports:
      - "8080:80"
    depends_on:
      ml-service:
        condition: service_healthy
    restart: unless-stopped
\end{lstlisting}

\newpage
%===================== ANNEXE 3 =====================%
\section*{Annexe 3.~Constantes et facteurs d'émission documentés}
\addcontentsline{toc}{section}{Annexe 3.~Constantes et facteurs d'émission documentés}

\addcontentsline{lot}{table}{Annexe 3.1~~~Facteurs d'émission utilisés pour la génération des données d'entraînement}

Le tableau annexe~3.1 recense les principaux facteurs d'émission physiques utilisés lors de la construction des jeux de données d'entraînement (script \texttt{prepare\_datasets.py}), avec leur source officielle.

{\raggedright \textbf{Tableau annexe 3.1~:}~Facteurs d'émission documentés}
\begin{longtable}[c]{
    | p{.42\textwidth}
    | p{.20\textwidth}
    | p{.30\textwidth} |
}
    \hline
    \textbf{Facteur} & \textbf{Valeur} & \textbf{Source} \\ \hline
    \endhead
    Facteur d'émission du réseau électrique (USA) & 0,42~kg CO2/kWh & EPA eGRID \\ \hline
    Chauffage au gaz naturel & 0,18~kg CO2e/kWh & DOE/EPA \\ \hline
    Chauffage au fioul & 0,27~kg CO2e/kWh & DOE/EPA \\ \hline
    Pompe à chaleur & 0,08~kg CO2e/kWh & DOE/EPA (COP~$\approx$~3) \\ \hline
    Intensité du gaspillage alimentaire & 2,5~kg CO2/kg & FAO / WRAP \\ \hline
    Débit d'une douche économe & 9,5~L/min & EPA WaterSense \\ \hline
    Équivalence arbre planté & 21~kg CO2/an & Estimation usuelle de séquestration \\ \hline
\end{longtable}
